\documentclass[10pt,a4paper]{amsart}
\usepackage{amsmath,amssymb,amsthm,graphicx}
% Retained minimal preamble for the Stacks Project chapter more-algebra.tex (More on Algebra).
% The chapter file pinned-original-source.tex carries only \input{preamble} on its first line, so no
% preamble text lives inside the chapter itself.  The Stacks Project uses one single shared preamble for
% every chapter; the snapshot stages it identically as sources/stacks-chow/preamble.tex and
% sources/stacks-etale-cohomology/preamble.tex.  The block below is copied mechanically, line for line,
% from sources/stacks-chow/preamble.tex lines 14-18 (the Xy-pic setup) and 154-217 (the native theorem
% declarations together with the mathematical macros).  Nothing was added, removed or redefined; the
% wrapper class and the amsmath/amssymb/amsthm/graphicx base packages are supplied by the assembler.
% Chapter: more-algebra.tex -- More on Algebra.
% For commutative diagrams we use Xy-pic
\usepackage[all]{xy}

% We use 2cell for 2-commutative diagrams.
\xyoption{2cell}

% Theorem environments.
%
\theoremstyle{plain}
\newtheorem{theorem}[subsection]{Theorem}
\newtheorem{proposition}[subsection]{Proposition}
\newtheorem{lemma}[subsection]{Lemma}

\theoremstyle{definition}
\newtheorem{definition}[subsection]{Definition}
\newtheorem{example}[subsection]{Example}
\newtheorem{exercise}[subsection]{Exercise}
\newtheorem{situation}[subsection]{Situation}

\theoremstyle{remark}
\newtheorem{remark}[subsection]{Remark}
\newtheorem{remarks}[subsection]{Remarks}

\numberwithin{equation}{subsection}

% Macros
%
\def\lim{\mathop{\mathrm{lim}}\nolimits}
\def\colim{\mathop{\mathrm{colim}}\nolimits}
\def\Spec{\mathop{\mathrm{Spec}}}
\def\Hom{\mathop{\mathrm{Hom}}\nolimits}
\def\Ext{\mathop{\mathrm{Ext}}\nolimits}
\def\SheafHom{\mathop{\mathcal{H}\!\mathit{om}}\nolimits}
\def\SheafExt{\mathop{\mathcal{E}\!\mathit{xt}}\nolimits}
\def\Sch{\mathit{Sch}}
\def\Mor{\mathop{\mathrm{Mor}}\nolimits}
\def\Ob{\mathop{\mathrm{Ob}}\nolimits}
\def\Sh{\mathop{\mathit{Sh}}\nolimits}
\def\NL{\mathop{N\!L}\nolimits}
\def\CH{\mathop{\mathrm{CH}}\nolimits}
\def\proetale{{pro\text{-}\acute{e}tale}}
\def\etale{{\acute{e}tale}}
\def\QCoh{\mathit{QCoh}}
\def\Ker{\mathop{\mathrm{Ker}}}
\def\Im{\mathop{\mathrm{Im}}}
\def\Coker{\mathop{\mathrm{Coker}}}
\def\Coim{\mathop{\mathrm{Coim}}}

% Boxtimes
%
\DeclareMathSymbol{\boxtimes}{\mathbin}{AMSa}{"02}

%
% Macros for moduli stacks/spaces
%
\def\QCohstack{\mathcal{QC}\!\mathit{oh}}
\def\Cohstack{\mathcal{C}\!\mathit{oh}}
\def\Spacesstack{\mathcal{S}\!\mathit{paces}}
\def\Quotfunctor{\mathrm{Quot}}
\def\Hilbfunctor{\mathrm{Hilb}}
\def\Curvesstack{\mathcal{C}\!\mathit{urves}}
\def\Polarizedstack{\mathcal{P}\!\mathit{olarized}}
\def\Complexesstack{\mathcal{C}\!\mathit{omplexes}}
% \Pic is the operator that assigns to X its picard group, usage \Pic(X)
% \Picardstack_{X/B} denotes the Picard stack of X over B
% \Picardfunctor_{X/B} denotes the Picard functor of X over B
\def\Pic{\mathop{\mathrm{Pic}}\nolimits}
\def\Picardstack{\mathcal{P}\!\mathit{ic}}
\def\Picardfunctor{\mathrm{Pic}}
\def\Deformationcategory{\mathcal{D}\!\mathit{ef}}

% --- appended 2026-10-12 (Chief): the project own definitions of the reference/slogan/history
% display environments, copied mechanically from the shared Stacks preamble lines 8-12
% (\usepackage{verbatim} plus the three \newenvironment lines, which the chapter preamble keeps
% verbatim from the project preamble). Nothing added, removed or redefined: these environments
% carry no mathematical content and the project itself typesets them as comments. Without these
% three lines a counted statement whose own text uses them fails with Environment slogan undefined.
\usepackage{verbatim}
\newenvironment{reference}{\comment}{\endcomment}
\newenvironment{slogan}{\comment}{\endcomment}
\newenvironment{history}{\comment}{\endcomment}
\usepackage{cleveref}
\crefname{originallemma}{original lemma}{original lemmas}
\Crefname{originallemma}{Original lemma}{Original lemmas}
\makeatletter
\newlabel{more-algebra-lemma-completed-direct-sum-flat}{{06LE}{source}{}{sourcepointer-1}{}}
\newlabel{more-algebra-lemma-completed-direct-sum-flat@cref}{{[originallemma][0][]06LE}{[1][0][]source}{}{}{}}
\makeatother
\crefname{originallemma}{original lemma}{original lemmas}
\Crefname{originallemma}{Original lemma}{Original lemmas}
\makeatletter
\newlabel{algebra-lemma-completion-flat}{{00MB}{source}{}{sourcepointer-2}{}}
\newlabel{algebra-lemma-completion-flat@cref}{{[originallemma][0][]00MB}{[1][0][]source}{}{}{}}
\makeatother
\crefname{originallemma}{original lemma}{original lemmas}
\Crefname{originallemma}{Original lemma}{Original lemmas}
\makeatletter
\newlabel{algebra-lemma-composition-flat}{{00HC}{source}{}{sourcepointer-3}{}}
\newlabel{algebra-lemma-composition-flat@cref}{{[originallemma][0][]00HC}{[1][0][]source}{}{}{}}
\makeatother
\crefname{originallemma}{original lemma}{original lemmas}
\Crefname{originallemma}{Original lemma}{Original lemmas}
\makeatletter
\newlabel{algebra-lemma-completion-complete}{{031C}{source}{}{sourcepointer-4}{}}
\newlabel{algebra-lemma-completion-complete@cref}{{[originallemma][0][]031C}{[1][0][]source}{}{}{}}
\makeatother
\crefname{originallemma}{original lemma}{original lemmas}
\Crefname{originallemma}{Original lemma}{Original lemmas}
\makeatletter
\newlabel{lemma-ui-completion-direct-sum-into-product}{{05BC}{source}{}{sourcepointer-5}{}}
\newlabel{lemma-ui-completion-direct-sum-into-product@cref}{{[originallemma][0][]05BC}{[1][0][]source}{}{}{}}
\makeatother
\crefname{originallemma}{original lemma}{original lemmas}
\Crefname{originallemma}{Original lemma}{Original lemmas}
\makeatletter
\newlabel{algebra-lemma-ui-flat-domain}{{058P}{source}{}{sourcepointer-6}{}}
\newlabel{algebra-lemma-ui-flat-domain@cref}{{[originallemma][0][]058P}{[1][0][]source}{}{}{}}
\makeatother
\crefname{originallemma}{original lemma}{original lemmas}
\Crefname{originallemma}{Original lemma}{Original lemmas}
\makeatletter
\newlabel{algebra-proposition-characterize-coherent}{{05CZ}{source}{}{sourcepointer-7}{}}
\newlabel{algebra-proposition-characterize-coherent@cref}{{[originallemma][0][]05CZ}{[1][0][]source}{}{}{}}
\makeatother
\crefname{originallemma}{original lemma}{original lemmas}
\Crefname{originallemma}{Original lemma}{Original lemmas}
\makeatletter
\newlabel{algebra-lemma-Noetherian-coherent}{{05CY}{source}{}{sourcepointer-8}{}}
\newlabel{algebra-lemma-Noetherian-coherent@cref}{{[originallemma][0][]05CY}{[1][0][]source}{}{}{}}
\makeatother
\title{The completion of a direct sum of copies of a Noetherian ring is flat}
\author{The Stacks Project Authors}
\date{stacks-project:more-algebra.tex:local-snapshot-20261012; GNU FDL 1.2 or later}
\begin{document}
\maketitle
\setcounter{subsection}{1}
\noindent\textit{Editorial note.} All delivered mathematical text is the author's own, copied from the pinned Stacks Project chapter file more-algebra.tex; no statement or proof step was rewritten, reordered or removed.\par\medskip
\noindent\textit{Editorial note.} The chapter file pinned-original-source.tex carries only \textbackslash{}input\{preamble\} on its first line, so no preamble text lives inside the chapter. The retained file sources/stacks-more-algebra/retained-preamble.tex is a mechanical, line-for-line copy of the project's single shared preamble as staged in the snapshot (Xy-pic setup plus theorem declarations and mathematical macros). Nothing was added, removed or redefined, and the excerpt serves only as preamble, never as a span coordinate.\par\medskip
\noindent\textit{Editorial note.} Both counted spans start at the \textbackslash{}begin line of their environment, so the statement and the author's proof are delivered whole; surrounding prose of the section is deliberately outside the counted spans and is carried as standard prerequisites instead.\par\medskip
\noindent\textit{Editorial note.} The citation fields (source\_id, source\_version, authors, source\_url, license.name, source\_title) are identical to those of every other card this lane has produced from this chapter, so the batch selector treats them as one source.\par\medskip
\noindent\textit{Allowed standard prerequisites (editorial).}\begin{itemize}
\item The completion map is flat, flatness is stable under composition, and completing a complete ring does not change it. Algebra, Lemma \textbackslash{}ref\{algebra-lemma-completion-flat\}, Algebra, Lemma \textbackslash{}ref\{algebra-lemma-composition-flat\}, and Algebra, Lemma \textbackslash{}ref\{algebra-lemma-completion-complete\}, cited at proof lines 6420-6431. Justify replacing the ring by its completion and rewriting the completed direct sum.
\item The comparison map from the completed direct sum into the product is universally injective. Lemma \textbackslash{}ref\{lemma-ui-completion-direct-sum-into-product\}, cited at proof lines 6432-6434. Reduces flatness of the completed direct sum to flatness of the product.
\item A universally injective map out of a module detects flatness. Algebra, Lemma \textbackslash{}ref\{algebra-lemma-ui-flat-domain\}, cited at proof lines 6435-6436. Carries flatness of the product back to the completed direct sum.
\item Products of copies of a coherent ring are flat, and Noetherian rings are coherent. Algebra, Proposition \textbackslash{}ref\{algebra-proposition-characterize-coherent\} and Algebra, Lemma \textbackslash{}ref\{algebra-lemma-Noetherian-coherent\}, cited at proof lines 6437-6440. Give flatness of the product of copies of the ring.
\end{itemize}\medskip
\ifdefined\hypertarget\hypertarget{sourcepointer-1}{}\fi
\noindent\textit{Source reference.} Original lemma 06LE: focused-seeds/AGP/sources/stacks-snapshot-20261011/tags.txt, source lines 8463--8463. The lemma's own canonical tag, quoted so that readers can resolve the label inside the delivered text; it points at the counted statement itself and carries no proof content.\par\medskip
\ifdefined\hypertarget\hypertarget{sourcepointer-2}{}\fi
\noindent\textit{Source reference.} Original lemma 00MB: focused-seeds/AGP/sources/stacks-snapshot-20261011/tags.txt, source lines 817--817. Navigation pointer to a statement the counted proof refers to but does not contain; nothing load-bearing is delegated to it.\par\medskip
\ifdefined\hypertarget\hypertarget{sourcepointer-3}{}\fi
\noindent\textit{Source reference.} Original lemma 00HC: focused-seeds/AGP/sources/stacks-snapshot-20261011/tags.txt, source lines 638--638. Navigation pointer to a statement the counted proof refers to but does not contain; nothing load-bearing is delegated to it.\par\medskip
\ifdefined\hypertarget\hypertarget{sourcepointer-4}{}\fi
\noindent\textit{Source reference.} Original lemma 031C: focused-seeds/AGP/sources/stacks-snapshot-20261011/tags.txt, source lines 3979--3979. Navigation pointer to a statement the counted proof refers to but does not contain; nothing load-bearing is delegated to it.\par\medskip
\ifdefined\hypertarget\hypertarget{sourcepointer-5}{}\fi
\noindent\textit{Source reference.} Original lemma 05BC: focused-seeds/AGP/sources/stacks-snapshot-20261011/tags.txt, source lines 6880--6880. Navigation pointer to a statement the counted proof refers to but does not contain; nothing load-bearing is delegated to it.\par\medskip
\ifdefined\hypertarget\hypertarget{sourcepointer-6}{}\fi
\noindent\textit{Source reference.} Original lemma 058P: focused-seeds/AGP/sources/stacks-snapshot-20261011/tags.txt, source lines 6787--6787. Navigation pointer to a statement the counted proof refers to but does not contain; nothing load-bearing is delegated to it.\par\medskip
\ifdefined\hypertarget\hypertarget{sourcepointer-7}{}\fi
\noindent\textit{Source reference.} Original lemma 05CZ: focused-seeds/AGP/sources/stacks-snapshot-20261011/tags.txt, source lines 6937--6937. Navigation pointer to a statement the counted proof refers to but does not contain; nothing load-bearing is delegated to it.\par\medskip
\ifdefined\hypertarget\hypertarget{sourcepointer-8}{}\fi
\noindent\textit{Source reference.} Original lemma 05CY: focused-seeds/AGP/sources/stacks-snapshot-20261011/tags.txt, source lines 6936--6936. Navigation pointer to a statement the counted proof refers to but does not contain; nothing load-bearing is delegated to it.\par\medskip
\begin{lemma}
\label{lemma-completed-direct-sum-flat}
Let $R$ be a ring. Let $I \subset R$ be an ideal. Let $A$ be a set.
Assume $R$ is Noetherian. The completion
$(\bigoplus\nolimits_{\alpha \in A} R)^\wedge$
is a flat $R$-module.
\end{lemma}

\begin{proof}
Denote $R^\wedge$ the completion of $R$ with respect to $I$. As
$R \to R^\wedge$ is flat by
Algebra, Lemma \ref{algebra-lemma-completion-flat}
it suffices to prove that
$(\bigoplus\nolimits_{\alpha \in A} R)^\wedge$ is a flat
$R^\wedge$-module (use
Algebra, Lemma \ref{algebra-lemma-composition-flat}).
Since
$$
(\bigoplus\nolimits_{\alpha \in A} R)^\wedge
=
(\bigoplus\nolimits_{\alpha \in A} R^\wedge)^\wedge
$$
we may replace $R$ by $R^\wedge$ and assume that $R$ is complete with
respect to $I$ (see
Algebra, Lemma \ref{algebra-lemma-completion-complete}).
In this case
Lemma \ref{lemma-ui-completion-direct-sum-into-product}
tells us the map
$(\bigoplus\nolimits_{\alpha \in A} R)^\wedge \to \prod_{\alpha \in A} R$
is universally injective. Thus, by
Algebra, Lemma \ref{algebra-lemma-ui-flat-domain}
it suffices to show that $\prod_{\alpha \in A} R$ is flat. By
Algebra, Proposition \ref{algebra-proposition-characterize-coherent}
(and
Algebra, Lemma \ref{algebra-lemma-Noetherian-coherent})
we see that $\prod_{\alpha \in A} R$ is flat.
\end{proof}

\end{document}
